\section*{Exercises on \S\arabic{exercisegroup}}
\phantomsection
\addcontentsline{toc}{section}{Exercises on \protect\S\arabic{exercisegroup}}

\begin{enumerate}
\def\labelenumi{\arabic{enumi})}
\tightlist
\item
  Let I be an arbitrary interval in R, let A be a set, and g a finite real function defined on \(I \times A\) such that for every \(\alpha \in A\) the map \(t \mapsto g(t, \alpha)\) is decreasing and \(\geqslant 0\) on I; suppose further that there is a number M independent of \(\alpha\) such that \(g(t, \alpha) \leqslant M\) on \(I \times A\) .
\end{enumerate}

\begin{enumerate}
\def\labelenumi{\alph{enumi})}
\item
  Show that if \(\mathbf{f}\) is a regulated function on I such that the integral \(\int_{\mathrm{I}} \mathbf{f}(t) dt\) converges, then the integral \(\int_{\mathrm{I}} \mathbf{f}(t) g(t, \alpha) dt\) is uniformly convergent for \(\alpha \in \mathrm{A}\) (use the second mean value theorem; cf.~II, p.~82, exerc. 16).
\item
  Suppose that I = {[}a, +∞{]} and also that \(g(t, \alpha)\) tends uniformly to 0 (for \(\alpha \in A\) ) when t tends to +∞. Show that if f is a regulated function on I and if there is a number k \textgreater{} 0 such that \(\left\|\int_{J} \mathbf{f}(t) dt\right\| \leqslant k\) for every compact interval J contained in I, then the integral \(\int_{I} \mathbf{f}(t) g(t, \alpha) dt\) is uniformly convergent for \(\alpha \in A\) (same method).
\item
  Suppose that \(I = [a, +\infty[\) with a \textgreater{} 0, \(A = [0, +\infty[,\) and \(g(t, \alpha) = \varphi(\alpha t)\) , where \(\varphi\) is a convex decreasing function and \(\geqslant 0\) on A, tending to 0 as x tends to \(+\infty\) and such that \(\varphi(0) = 1\) . Suppose further that \(f = h''\) , where h is a twice differentiable function on \([a, +\infty[,\) and, together with \(h'\) , is bounded on this interval. The integral \(\int_{a}^{+\infty} f(t) \varphi(\alpha t) dt\) is then convergent for \(\alpha > 0\) ; show that when \(\alpha\) tends to 0, it tends to \(-\mathbf{h}'(a)\) (integrate by parts and use the second mean value theorem).
\end{enumerate}

* \(d\) ) Take \(\varphi (x) = e^{-cx}\) , where \(c > 0\) , take \(a = 1\) , and for \(f\) take the complex function \(t\mapsto e^{i\log t} / t\) , which is the derivative of a bounded function on I; show that for a suitable choice of \(c\) the integral \(\int_{1}^{+\infty}e^{-\alpha ct}\frac{e^{i\log t}}{t} dt\) , which is absolutely convergent for \(\alpha >0\) , does not tend to any limit when \(\alpha\) tends to 0 (after integrating by parts, make the change of variables \(\alpha t = u\) ; use Laplace transform theory to see that \(\int_0^{+\infty}e^{-cu + i\log u}du\) does not vanish for certain \(c > 0\) .)*

\begin{enumerate}
\def\labelenumi{\arabic{enumi})}
\setcounter{enumi}{1}
\tightlist
\item
  Let \(f\) and \(g\) be two real regulated functions on a compact interval \([a, b]\) , such that \(f\) is decreasing on \([a, b]\) , and \(0 \leqslant g(t) \leqslant 1\) . If one puts \(\lambda = \int_{a}^{b} g(t) dt\) , show that one has
\end{enumerate}

\[
\int_ {b - \lambda} ^ {b} f (t) d t <   \int_ {a} ^ {b} f (t) g (t) d t <   \int_ {a} ^ {a + \lambda} f (t) d t
\]

except when f is constant, or g is equal to 0 (resp. 1) at all the points where it is continuous (in which case all three terms are equal). (Vary one of the limits of integration in the integral \(\int_{x}^{y} f(t)g(t)dt\) .)

\begin{enumerate}
\def\labelenumi{\arabic{enumi})}
\setcounter{enumi}{2}
\tightlist
\item
  Let \(f(x, \alpha) = 1 / \sqrt{1 - 2\alpha x + \alpha^2}\) for \(-1 < x < +1\) and \(\alpha \in \mathbf{R}\) ; show that the function \(g(\alpha) = \int_{-1}^{+1} dx / \sqrt{1 - 2\alpha x + \alpha^2}\) is continuous on \(\mathbf{R}\) , but has no derivative at \(\alpha = 1\) and \(\alpha = -1\) ; show that \(f_{\alpha}'(x, \alpha)\) exists for all \(\alpha \in \mathbf{R}\) and for \(x \in I = ] - 1, +1[\) , and is continuous on \(I \times R\) , and that the integral \(\int_{-1}^{+1} f_{\alpha}'(x, \alpha) dx\) exists for all \(\alpha \in R\) , but verify that this integral is not uniformly convergent on a neighbourhood of the point \(\alpha = 1\) or of the point \(\alpha = -1\) .
\end{enumerate}

¶ 4) Let I be an interval of \(\mathbf{R}\) , A a neighbourhood of a point \(\alpha_0\) in the field \(\mathbf{R}\) (resp. the field \(\mathbf{C}\) ), and \(\mathbf{f}\) a continuous map of \(\mathrm{I} \times \mathrm{A}\) into a complete normed space \(\mathrm{E}\) over \(\mathbf{R}\) , such that \(\mathbf{f}_{\alpha}'(x, \alpha)\) exists and is continuous on \(\mathrm{I} \times \mathrm{A}\) . Let \(a(\alpha)\) , \(b(\alpha)\) be two continuous functions defined on A, with values in I, such that one has identically \(\mathbf{f}(a(\alpha), \alpha) = \mathbf{f}(b(\alpha), \alpha) = 0\) on A. Show that the function \(\mathbf{g}(\alpha) = \int_{a(\alpha)}^{b(\alpha)} \mathbf{f}(t, \alpha) dt\) admits a derivative equal to \(\int_{a(\alpha_0)}^{b(\alpha_0)} \mathbf{f}_{\alpha}'(t, \alpha_0) dt\) at the point \(\alpha_0\) , even if \(a\) and \(b\) are not differentiable at the point \(\alpha_0\) (let M be the supremum of \(\left\| \mathbf{f}_{\alpha}'(x, \alpha) \right\|\) on a compact neighbourhood of \((b(\alpha_0), \alpha_0)\) ; note, applying Bolzano's theorem to \(b(\alpha)\) , that for every \(x\) belonging to the interval with endpoints \(b(\alpha_0)\) and \(b(\alpha)\) , one has, for \(\alpha\) sufficiently close to \(\alpha_0\) , that \(\| \mathbf{f}(x, \alpha) \| \leqslant M |\alpha - \alpha_0|\) ).

\begin{enumerate}
\def\labelenumi{\arabic{enumi})}
\setcounter{enumi}{4}
\tightlist
\item
  Let \(\mathbf{f}\) be a continuous vector function on a compact interval \(I = [0, a]\) . Show that if, at the point \(\alpha_0 \in I\) , there is an \(\varepsilon > 0\) such that \((\mathbf{f}(x) - \mathbf{f}(\alpha_0)) / |x - \alpha_0|^{1/2 + \varepsilon}\) remains bounded as \(x\) tends to \(\alpha_0\) , then the function \(\mathbf{g}(\alpha) = \int_0^\alpha \mathbf{f}(x) dx / \sqrt{\alpha - x}\) admits a derivative equal to
\end{enumerate}

\[
\frac {1}{\sqrt {\alpha_ {0}}} \mathbf {f} (\alpha_ {0}) - \frac {1}{2} \int_ {0} ^ {\alpha_ {0}} \frac {\mathbf {f} (x) - \mathbf {f} (\alpha_ {0})}{(\alpha_ {0} - x) ^ {3 / 2}} d x
\]

at the point \(\alpha_0\) for \(\alpha_0 > 0\) , and to 0 for \(\alpha = 0\) .

When \(\mathbf{f}\) is the real function \(\sqrt{\alpha_0 - x}\) , show that the function \(\mathbf{g}\) has an infinite derivative at the point \(\alpha_0\) .

¶ 6) Let \(\mathrm{I} = [a, b]\) , \(\mathrm{A} = [c, d]\) be two compact intervals on \(\mathbf{R}\) ; let \(\mathbf{f}\) be a function defined on \(\mathrm{I} \times \mathrm{A}\) , with values in a complete normed space \(\mathrm{E}\) over \(\mathbf{R}\) , such that for all \(\alpha \in \mathrm{A}\) , the map \(t \mapsto \mathbf{f}(t, \alpha)\) is regulated on \(\mathrm{I}\) , that \(\mathbf{f}\) is bounded on \(\mathrm{I} \times \mathrm{A}\) , and that the set \(\mathrm{D}\) of points of discontinuity of \(\mathbf{f}\) in \(\mathrm{I} \times \mathrm{A}\) is met in a finite number of points by each line \(x = x_0\) and each line \(\alpha = \alpha_0\) ( \(x_0 \in \mathrm{I}\) , \(\alpha_0 \in \mathrm{A}\) ).

\begin{enumerate}
\def\labelenumi{\alph{enumi})}
\item
  Show that the function \(\mathbf{g}(\alpha) = \int_{a}^{b} \mathbf{f}(t, \alpha) dt\) is continuous on A (given \(\alpha_0 \in \mathrm{A}\) and \(\varepsilon > 0\) , show that there is a neighbourhood V of \(\alpha_0\) and a finite number of intervals \(J_k\) contained in I with the sum of their lengths \(\leqslant \varepsilon\) , such that, if J denotes the complement in I of \(\bigcup_{k} J_k\) , then \(\mathbf{f}\) is continuous on \(\mathrm{J} \times \mathrm{V}\) ).
\item
  Show that the formula for interchanging the order of integration (II, p.~77, formula (15)) is still valid (same method as in a)).
\end{enumerate}

\begin{enumerate}
\def\labelenumi{\arabic{enumi})}
\setcounter{enumi}{6}
\tightlist
\item
  Let \(f\) be a real function defined and having a continuous derivative on the interval \(]0, +\infty[\) , and such that
\end{enumerate}

\[
\lim _ {x \to 0} f (x) = \lim _ {x \to + \infty} f (x) = 0.
\]

The integral \(\int_0^{+\infty}f'(\alpha t)dt\) is defined and continuous on every bounded interval \(]0,a]\) ; show that the integral \(\int_0^{+\infty}dt\int_0^a f'(\alpha t)d\alpha\) may either not exist or may be different from

\[
\int_ {0} ^ {a} d \alpha \int_ {0} ^ {+ \infty} f ^ {\prime} (\alpha t) d t.
\]

¶ 8) Let \(I\) and \(J\) be two arbitrary intervals in \(\mathbf{R}\) , and \(\mathbf{f}\) a function defined and continuous on \(I \times J\) , with values in a complete normed space \(E\) over \(\mathbf{R}\) . Suppose that:

\(1^{\circ}\) the integral \(\int_{\mathrm{I}}\mathbf{f}(x,y)dx\) is uniformly convergent when \(y\) runs through an arbitrary compact interval contained in J;

\(2^{\circ}\) the integral \(\int_{\mathrm{J}}\mathbf{f}(x,y)dy\) is uniformly convergent when \(x\) runs through an arbitrary compact interval contained in I;

\(3^{\circ}\) if, for every compact interval \(\mathrm{H}\) contained in I, one puts \(\mathbf{u}_{\mathrm{H}}(y) = \int_{\mathrm{H}} \mathbf{f}(x, y) dx\) , the integral \(\int_{\mathrm{J}} \mathbf{u}_{\mathrm{H}}(y) dy\) is uniformly convergent for \(\mathrm{H} \in \mathfrak{K}(\mathrm{I})\) (the right directed ordered set of compact intervals contained in I).

Under these conditions, show that the integrals \(\int_{\mathrm{I}} dx \int_{\mathrm{J}} \mathbf{f}(x, y) dy\) and \(\int_{\mathrm{J}} dy \int_{\mathrm{I}} \mathbf{f}(x, y) dx\) exist and are equal.

* 9) Deduce from exerc. 8 that the integrals

\[
\int_ {0} ^ {\infty} d x \int_ {0} ^ {\infty} e ^ {- y x ^ {2}} \sin y d y \quad \text { and } \quad \int_ {0} ^ {\infty} d y \int_ {0} ^ {\infty} e ^ {- y x ^ {2}} \sin y d x
\]

exist and are equal.*

\begin{enumerate}
\def\labelenumi{\arabic{enumi})}
\setcounter{enumi}{9}
\tightlist
\item
  \begin{enumerate}
  \def\labelenumii{\alph{enumii})}
  \tightlist
  \item
    If \(h, u, v'\) are primitives of real regulated functions on an open interval \(]a, b[\) of \(\mathbf{R}\) , if \(v\) is a primitive of \(v'\) , and if \(v(x) > 0\) on this interval, then one has the Redheffer identity
  \end{enumerate}
\end{enumerate}

\[
h u ^ {\prime 2} = - h \mathrm{D} \left(\frac {h v ^ {\prime}}{v}\right) + h v ^ {2} \left(\mathrm{D} \left(\frac {u}{v}\right)\right) ^ {2} + \mathrm{D} \left(\frac {u ^ {2} h v ^ {\prime}}{v}\right)
\]

at those points of {]}a, b{[} where the derivatives are defined.

* b) Let \(v\) , \(w\) be two functions \(> 0\) on \(]a\) , \(b[\) , which are primitives of regulated functions \(v' > 0\) , \(w' \leqslant 0\) . Deduce from \(a)\) that for every function \(u\) which is a primitive of a regulated function \(u'\) on \(]a\) , \(b[\) and such that \(\liminf_{x \to a, x > a} u(x) = 0\) , the hypothesis that the integral \(\int_{a}^{b} \frac{w(x)}{v'(x)} u'^2(x) dx\) is convergent implies that the integral \(\int_{a}^{b} \frac{w'(x)}{v(x)} u^2(x) dx\) is convergent and that

\[
\limsup _ {x \to b, x <   b} u ^ {2} (x) w (x) / v (x)
\]

is finite, and also the inequality

\[
\int_ {a} ^ {b} \frac {w (x)}{v ^ {\prime} (x)} u ^ {\prime 2} (x) d x \geqslant - \int_ {a} ^ {b} \frac {w ^ {\prime} (x)}{v (x)} u ^ {2} (x) d x + \operatorname * {l i m s u p} _ {x \to b, x <   b} \frac {u ^ {2} (x) w (x)}{v (x)}.\tag{*}
\]

(Putting \(h = w / v'\) , first observe that \(\int_{c}^{x} dt / h(t) \leqslant v(x) / w(x)\) for \(a < c < x < b\) and remark that by the Cauchy-Schwarz inequality

\[
(u (x) - u (c)) ^ {2} \leqslant \left(\int_ {c} ^ {x} h (t) u ^ {\prime 2} (t) d t\right) \left(\int_ {c} ^ {x} \frac {d t}{h (t)}\right),\tag{**}
\]

to deduce that

\[
\liminf _ {x \to a, x > a} u ^ {2} (x) w (x) / v (x) = 0,
\]

and then integrate the Redheffer identity.)

\begin{enumerate}
\def\labelenumi{\alph{enumi})}
\setcounter{enumi}{2}
\tightlist
\item
  Deduce from (*) that if u is the primitive of a regulated function on {]}0, 1{]}, if
\end{enumerate}

\[
\liminf _ {x \to 0, x > 0} u (x) = 0
\]

and if the integral \(\int_0^1 u'^2 (t)dt\) converges, then so does \(\int_0^1 (u(t) / t)^2 dt\) and for all \(\alpha > 0\) one has the inequality

\[
\int_ {0} ^ {1} u ^ {\prime 2} (t) d t \geqslant \alpha (1 - \alpha) \int_ {0} ^ {1} \left(\frac {u (t)}{t}\right) ^ {2} d t + \alpha u ^ {2} (1).
\]

\begin{enumerate}
\def\labelenumi{\alph{enumi})}
\setcounter{enumi}{3}
\tightlist
\item
  Let \(\alpha\) be a number \textgreater0, let K be a function \textgreater0, differentiable and decreasing on \(]0, +\infty[\) and such that \(\lim_{x \to +\infty} K(x) = 0\) . If u is the primitive of a regulated function on \(]0, +\infty[,\) such that \(\liminf_{x \to 0, x > 0} u(x) = 0\) , and if the integral \(\int_{0}^{+\infty} x^{1-\alpha} K(x) u'^{2}(x) dx\) converges, then the function \(x^{-\alpha} K(x) u^{2}(x)\) tends to 0 as x tends to 0 or to \(+\infty\) , the integral
\end{enumerate}

\[
\int_ {0} ^ {+ \infty} x ^ {- \alpha} \mathrm{K} ^ {\prime} (x) u ^ {2} (x) d x
\]

converges, and one has

\[
\int_ {0} ^ {+ \infty} x ^ {1 - \alpha} \mathrm{K} (x) u ^ {\prime 2} (x) d x \geqslant - \alpha \int_ {0} ^ {+ \infty} x ^ {- \alpha} \mathrm{K} ^ {\prime} (x) u ^ {2} (x) d x
\]

(take \(v(x) = x^{\alpha}\) , \(h(x) = x^{1 - \alpha}\mathrm{K}(x)\) in the Redheffer identity).

In particular, for \(\alpha = \frac{1}{2}\) and \(\mathrm{K}(x) = x^{-1 / 2}\) , one has

\[
4 \int_ {0} ^ {+ \infty} u ^ {\prime 2} (x) d x \geqslant \int_ {0} ^ {+ \infty} \left(\frac {u (x)}{x}\right) ^ {2} d x
\]

(Hardy-Littlewood inequality).

\begin{enumerate}
\def\labelenumi{\alph{enumi})}
\setcounter{enumi}{4}
\tightlist
\item
  Let \(\alpha \geqslant -1\) , let \(K\) be a function \(\geqslant 0\) , differentiable and increasing on \([0, +\infty[\) . Suppose that the integrals
\end{enumerate}

\[
\int_ {0} ^ {+ \infty} x ^ {- \alpha} \mathrm{K} (x) u ^ {\prime 2} (x) d x \quad \text { and } \quad \int_ {0} ^ {+ \infty} x ^ {\alpha} \mathrm{K} (x) u ^ {2} (x) d x
\]

converge. Then \(\mathrm{K}(x)u^{2}(x)\) tends to 0 as x tends to \(+\infty\) , and one has

\[
\int_ {0} ^ {+ \infty} \mathrm{K} ^ {\prime} (x) u ^ {2} (x) d x \leqslant 2 \left(\int_ {0} ^ {+ \infty} x ^ {- \alpha} \mathrm{K} (x) u ^ {\prime 2} (x) d x\right) ^ {1 / 2} \left(\int_ {0} ^ {+ \infty} x ^ {\alpha} \mathrm{K} (x) u ^ {2} (x) d x\right) ^ {1 / 2}
\]

(take \(v(x) = \exp(-cx^{\alpha})\) in the Redheffer identity, with a suitable constant \(c\) ).

In particular, if a and b are constants such that \(b + 1 \geqslant a\) and \(a + b \geqslant 0\) , one has

\[
(a + b) \int_ {0} ^ {+ \infty} x ^ {a + b - 1} u ^ {2} (x) d x \leqslant 2 \left(\int_ {0} ^ {+ \infty} x ^ {2 a} u ^ {\prime 2} (x) d x\right) ^ {1 / 2} \left(\int_ {0} ^ {+ \infty} x ^ {2 b} u ^ {2} (x) d x\right) ^ {1 / 2}
\]

if the integrals on the right-hand side converge (generalized H. Weyl inequality).

\(f)\) If \(0 < \alpha < 2\) and \(u\) is a primitive of a regulated function on \([0, \alpha]\) and \(u(0) = 0\) , one has

\[
\int_ {0} ^ {\alpha} (\alpha - x) ^ {\alpha - 1} x ^ {1 - \alpha} u ^ {\prime} (x) \left(u ^ {\prime} (x) - 2 u (x)\right) d x \geqslant 0
\]

(generalized Opial inequality). (Apply (*) suitably, replacing \(u(x)\) by \(e^{-x}u(x)\) .)

\(g\) ) If \(u\) is a primitive of a regulated function on \([0,b]\) and \(u(0) = 0\) , then

\[
\int_ {0} ^ {b} e ^ {- 2 x} u ^ {\prime 2} (x) d x \geqslant \int_ {0} ^ {b} e ^ {- 2 x} u (x) u ^ {\prime} (x) d x + \frac {1}{2} u ^ {2} (b) e ^ {- 2 b}
\]

(Hlawka's inequality) (same method as in \(f\) )).

\begin{enumerate}
\def\labelenumi{\alph{enumi})}
\setcounter{enumi}{7}
\tightlist
\item
  If u is the primitive of a regulated function on R then, for all \(t \in R\) ,
\end{enumerate}

\[
| u (t) | \leqslant \left(\int_ {- \infty} ^ {+ \infty} u ^ {\prime 2} (x) d x\right) ^ {1 / 4} \left(\int_ {- \infty} ^ {+ \infty} u ^ {2} (x) d x\right) ^ {1 / 4}
\]

if the two integrals on the right-hand side converge (consider the two intervals \(]-\infty,t]\) and \([t,+\infty[\) ; take \(v(x)=e^{\alpha x}\) on the two intervals, \(h(x)=1/\alpha\) on the first interval and \(h(x)=-1/\alpha\) on the second, then choose \(\alpha>0\) suitably). \(^{*}\)

